\documentclass[11pt]{article}
\usepackage[T1]{fontenc}
\usepackage{lmodern}
\usepackage{booktabs}
\usepackage{hyperref}

\title{Sourceright: Evidence-Bounded Reference Verification for Scholarly Workflows}
\author{Sourceright contributors}
\date{Technical preview --- version 0.1.20}

\begin{document}
\maketitle

\begin{abstract}
Scholarly writing workflows often need to distinguish a citation that is
syntactically present from a reference that has been reconciled against
external evidence. Sourceright is an open-source technical preview for that
workflow. It uses canonical CSL JSON for clean bibliographic records, a
separate verification sidecar for provider evidence and conflicts, and a
derived review queue for human action. This paper describes the architecture,
claim boundaries, deterministic fixtures, optional extraction backends, and
integration contracts for citation managers and journal or preprint workflows.
The implementation is intentionally evidence-bounded: provider observations
do not silently overwrite canonical data, and a warning or match is not a
claim that a source or proposition is true. We report the current engineering
surface and a reproducible evaluation plan rather than presenting unsupported
accuracy or production-readiness claims.
\end{abstract}

\section{Motivation}

Reference checking is commonly treated as a formatting task. In practice, a
workflow may need to preserve the author's record, compare it with provider
observations, identify missing or conflicting fields, and route uncertain cases
to review. These concerns become more difficult when the same manuscript moves
through a citation manager, a journal submission system, a preprint server, or
an institutional repository.

Sourceright addresses this as an audit-oriented workflow rather than as an
automatic truth judge. Its current repository contains Rust CLI and library
surfaces, fixture-backed provider adapters, report generation, and optional
read-only integration contracts. The software is a technical preview; its
fixture benchmarks are regression scaffolds, not externally comparable
performance measurements.

\section{Data and boundary model}

The core data boundary has three distinct layers:

\begin{enumerate}
  \item \texttt{references.csl.json} is canonical clean bibliographic data.
  \item \texttt{references.verification.json} stores provider observations,
        provenance, confidence inputs, conflicts, and review state.
  \item \texttt{review-queue.jsonl} is derived operational work for a person
        or downstream workflow.
\end{enumerate}

This separation is a safety property. An external provider can suggest a
candidate or expose a conflict without silently changing the author's
canonical record. Legal citations remain a separate model and are not folded
into academic CSL. Claim/source/provenance features may link evidence to
claims, but they do not assert the truth of a claim.

\section{Workflow architecture}

The normal path is intentionally inspectable:

\begin{enumerate}
  \item ingest a manuscript, citation file, or provider fixture;
  \item normalize records into the canonical CSL boundary;
  \item query selected providers or optional extraction services;
  \item record observations and provenance in the verification sidecar;
  \item emit deterministic reports and review items; and
  \item leave any mutation to an explicitly approved, audited workflow.
\end{enumerate}

The optional scholarly-document extraction boundary is backend-neutral. The
repository includes an explicitly selected GROBID REST adapter and native
Rust extraction experiments behind neutral schemas. The neutral layer is
GROBID-inspired, not a GROBID fork or a claim of API compatibility. GROBID is
an optional service boundary; it is not required for the core reference
workflow.

\section{Reproducibility and quality controls}

The repository keeps self-authored and license-attributed fixtures under
version control. Policy tests check schema shape, provenance fields, path
grounding, and separation from canonical and verification files. Workflow
controls pin GitHub Actions, use least-privilege permissions, disable
persisted checkout credentials, enforce concurrency, and bound job duration.
Release workflows generate checksums, attestations, and a Cargo dependency
metadata SBOM artifact.

The current evaluation strategy separates mandatory, network-free Rust checks
from optional integration lanes. Fixture-backed tests establish deterministic
regressions. Scheduled lanes may exercise coverage, fuzzing, mutation,
stress, and benchmark scaffolds, but a local fixture result is not presented
as state-of-the-art or population-level accuracy.

\begin{table}[ht]
\centering
\caption{Evidence levels used by the project.}
\begin{tabular}{ll}
\toprule
Level & Meaning \\
\midrule
Contracted & Interfaces, schemas, and policy boundaries are documented. \\
Fixture-backed & Deterministic local inputs and expected outputs pass. \\
Submitted & An external issue or submission URL exists. \\
Publicly accepted & An external maintainer or registry provides acceptance evidence. \\
\bottomrule
\end{tabular}
\end{table}

\section{Integration surfaces}

Sourceright can be used as a local CLI, a read-only MCP surface, or a thin
adapter in citation-manager, journal, and repository workflows. The current
integration posture is deliberately conservative. OJS, Janeway, arXiv
submit-ce, and other journal or repository platforms have fixture-backed or
contracted paths, but these should not be described as supported integrations
until a disposable instance, platform-specific smoke, permissions review, and
upstream contribution path have been verified.

For arXiv-style source submissions, source hygiene is part of the proposed
workflow. A submission bundle should contain only the files needed to compile
the paper, include the bibliography inputs required by the selected TeX
processor, and exclude credentials, repository archives, hidden files, and
unrelated working material. Sourceright can screen citation and provenance
signals before submission, but it does not submit to arXiv or mutate arXiv
state.

\section{Limitations and future work}

The current release does not establish universal citation accuracy, complete
PDF understanding, semantic claim verification, or institutional production
readiness. Live provider availability, parser disagreement, multilingual
coverage, and document-layout variation remain sources of uncertainty. Future
work includes differential conversion checks against carefully licensed
external tools, disposable-instance trials for open journal and repository
platforms, stronger release evidence, and maintainer-reviewed upstream
integration hooks.

\section{Conclusion}

Sourceright provides a conservative engineering pattern for reference
verification: preserve canonical data, isolate observations, expose
uncertainty, and make review work explicit. The repository's current evidence
supports a technical-preview and pilot-oriented description. It does not
support claims of final verification, benchmark superiority, or external
platform acceptance.

\bibliographystyle{plain}
\bibliography{references}
\end{document}
